% =====================================================================================================
% Conference paper skeleton: cooperative HARQ for NR sidelink platoon groupcast
% IEEE conference format (IEEEtran), 6 pages + references is the usual budget (VTC / ICC / Globecom).
%
% How to use:
%   - Every \todo{...} is a gap to fill. Every % NOTE: line is guidance; delete it when done.
%   - Numbers in [brackets] inside comments are results from the current campaigns that you may use;
%     re-check each one against its notebook before it goes into the text.
%   - Figures point to harq_helper_study/figures/ (see \graphicspath); the figure files exist already.
%   - Build: latexmk -pdf main.tex
% =====================================================================================================
\documentclass[conference]{IEEEtran}
\IEEEoverridecommandlockouts
\usepackage{cite}
\usepackage{amsmath,amssymb,amsfonts}
\usepackage{graphicx}
\usepackage{booktabs}
\usepackage{array}
\usepackage{xcolor}
\usepackage[hidelinks]{hyperref}

\graphicspath{{../scratch_overlay/platoon_harq_robustness_eval/harq_helper_study/figures/}}
\newcommand{\todo}[1]{\textcolor{red}{[TODO: #1]}}
% Scheme names: \scheme{<key>} prints the full name in running text, \Scheme{<key>} the capitalised
% form (sentence start, table rows). Keys are the simulator arm names; an unknown key prints a red TODO.
% Change a name in one place here.
\makeatletter
\newcommand{\defscheme}[3]{\@namedef{scheme@#1}{#2}\@namedef{Scheme@#1}{#3}}
\newcommand{\scheme}[1]{\@ifundefined{scheme@#1}{\todo{unknown scheme #1}}{\@nameuse{scheme@#1}}}
\newcommand{\Scheme}[1]{\@ifundefined{Scheme@#1}{\todo{unknown scheme #1}}{\@nameuse{Scheme@#1}}}
\makeatother
\defscheme{R3}{source Option~2 HARQ}{Source Option~2 HARQ}
\defscheme{R4+C}{blind helper relaying}{Blind helper relaying}
\defscheme{R4+C+fb}{blind helper relaying with fallback}{Blind helper relaying with fallback}
\defscheme{R4+C+fbs}{blind helper relaying with fallback and spaced reservations}%
          {Blind helper relaying with fallback and spaced reservations}
\defscheme{R6}{cooperative HARQ}{Cooperative HARQ}

\begin{document}

% -----------------------------------------------------------------------------------------------------
\title{\todo{Title}\\
% NOTE: name the problem and the answer, e.g. "Cooperative HARQ Relaying for NR Sidelink Platoon
% Groupcast: When Does a Helper Pay Off?"
}

\author{
\IEEEauthorblockN{\todo{Author 1}, \todo{Author 2}}
\IEEEauthorblockA{\todo{Affiliation, city, country}\\
\todo{email}}
}

\maketitle

% -----------------------------------------------------------------------------------------------------
\begin{abstract}
% NOTE: 150-200 words. Numbers checked on 2026-11-23 against campaign v3 and the helper-miss sweep:
%   - vs Option 2 (R3) under blockage: up to 18.7 pp (Roundabout E, aperiodic, no background;
%     17.8 pp on Roundabout W periodic). "up to 18 pp" is fine, "up to 19 pp" is more exact.
%   - fewer resources than Option 2: R6 uses 2.4-11.5 % fewer transmissions per packet in every cell
%     (max 11.5 %, Roundabout E periodic, no background): "up to 12 %" OK, or "up to 11.5 %".
%   - vs blind helper when the helper misses its first copy (p = 0.7, periodic, no background):
%     helper drops the payload: +17.3 pp highway, +24.6 pp Roundabout W, +25.5 pp Roundabout E;
%     helper drops the SCI: +13.0 / +14.3 / +14.4 pp. The draft's "17 pp" matches only the highway
%     payload case (or the 60-vehicle highway aperiodic load case without forced miss: +17.3 pp).
%   - resources vs blind helper: R6 uses FEWER at p = 0 (e.g. 1.75 vs 2.86 tx/pkt highway periodic) and
%     MORE only when takeovers fire often (helper drops the payload at high p on the roundabouts: 2.79 vs
%     1.58 at p = 0.7) or at high load with aperiodic traffic (2.68 vs 2.39, Roundabout W, 60 veh).
3GPP sidelink (SL) groupcast hybrid automatic repeat request (HARQ) provides feedback-controlled
temporal diversity by triggering retransmissions from the source. However, these retransmissions remain
vulnerable to persistent blockage. Blind helper relaying adds spatial diversity without group feedback,
but suffers when the helper misses the first transmission. We propose a cooperative groupcast HARQ that
delegates selected retransmission reservations to a helper, while receivers combine source and helper
copies in one HARQ process. A helper negative acknowledgment (NACK) triggers a source takeover, and a
helper that fails to retransmit triggers a fallback to source transmission. Using ns-3 simulations, we
show a gain of up to \todo{17\,pp (highway) / 25\,pp (roundabouts): choose, see NOTE} percentage
points (pp) against blind helper relaying when the helper misses the first transmission, and up to
\todo{18}\,pp against HARQ Option~2 in blockage scenarios. The proposed method uses up to \todo{12}\,\% fewer transmissions than Option~2, and
\todo{fewer than the blind helper without helper failures, equal or more when takeover and fallback
are triggered often}.
\end{abstract}

\begin{IEEEkeywords}
NR V2X, sidelink, groupcast, HARQ, cooperative retransmissions, Mode~2.
\end{IEEEkeywords}

% =====================================================================================================
\section{Introduction}\label{sec:intro}
% NOTE: about 1 column. Paragraph plan:
%   P1  Platooning and why reliability of groupcast matters (control loop, PDB = 100 ms periodic).
%   P2  NR sidelink groupcast HARQ: Option 1 (NACK-only) and Option 2 (ACK/NACK) and their limits:
%       a member that is blocked from the leader stays blocked for every retransmission.
%   P3  Relaying by a platoon member (helper): blind relaying is simple but adds load and collisions;
%       feedback-driven relaying is efficient but depends on PSFCH and on the helper itself.
%   P4  Gap: \todo{what prior work does not cover (system-level, standard-compliant, under load,
%       helper failure)}.
%   P5  Contributions (bullet list below).
%   P6  Paper organisation (one sentence).
\todo{Introduction}

\noindent Our contributions are:
\begin{itemize}
  \item \todo{A \scheme{R6} protocol for NR sidelink groupcast compatible with Rel-16/17
        signalling: helper named in SCI-2, relaying on NACK, leader takeover and fallback.}
  \item \todo{A system-level evaluation in ns-3 across three road scenarios, two traffic types and five
        channel-busy-ratio (CBR) levels, against leader-only HARQ and blind helper relaying.}
  \item \todo{A robustness study: forced helper failures (payload vs.\ SCI loss) and PSFCH errors.}
  \item \todo{A design lesson on resource selection: the PSFCH minimum time gap is what lets the leader
        recover from a silent helper, also without feedback (spaced blind arm).}
  % NOTE: keep 3-4 bullets; drop the last one if space is short.
\end{itemize}

% =====================================================================================================
\section{Background and Related Work}\label{sec:related}
\subsection{NR Sidelink Groupcast HARQ}
% NOTE: still missing from this paragraph: the PSFCH minimum time gap between a PSSCH and its feedback,
% and that it constrains how close the reserved retransmission slots can be when feedback is enabled
% (TS 38.214 8.1.4) -- Sec. VI uses this.
In Mode~2, the transmitting MAC selects SL resources from physical-layer (PHY) sensing results, as
specified in TS~38.214 and TS~38.321~\cite{3gpp_ts38214,3gpp_ts38321}. The selected set can contain the
initial transmission and later retransmission opportunities (up to three per transport block (TB),
depending on the pool configuration). First-stage sidelink control information (SCI) format~1-A
announces the reservation to nearby UEs. Second-stage SCI carries the source and destination identifiers
(IDs), HARQ process ID, new data indicator (NDI), and redundancy version (RV)~\cite{3gpp_ts38212}. These
fields let a receiver associate successive copies with one TB. The physical sidelink feedback channel
(PSFCH) carries HARQ feedback: groupcast Option~1 uses NACK-only feedback, whereas Option~2 assigns
receiver-specific ACK/NACK resources~\cite{3gpp_ts38213}.

\subsection{Cooperative Relaying for V2X}
% NOTE: every citation below must be checked (authors, year, venue, and that the paper says what we
% claim). kuo2021reliable checked in Zotero (2026-11-23): the source "reserves the retransmission
% resources and assigns the UEs in the same group to retransmit the TB"; at most two retransmission
% resources per initial transmission; helper selected from CAM positions; blind retransmission.
% leibel2020harq checked in Zotero (2026-11-23): Sec. IV, the transmitter selects the forwarding UE from the
%   HARQ feedback (decoded it, furthest toward the out-of-range members); multi-hop forwarding as in their [13];
%   no combining with the source's copies is described. Mainly a mixed NACK-only / ACK-NACK feedback paper.
% aoki2026enhancing (not in Zotero; read on mdpi.com 2026-11-23): Takahashi, Aoki, Tang, Sensors 2026.
%   BROADCAST at NLOS intersections (not platoon groupcast), Scenargie. The relay uses the source's reserved
%   retransmission resources; an explicit ACK confirms the relay decoded it (ideal error-free ACK assumed);
%   "timeout-based fallback" is named as possible but "not evaluated". Delivery modelled as independent
%   per-link trials (their Eq. 2). Extends their APCC 2024 paper.
% kim2019efficient (not in Zotero; read on IEEE Xplore 2026-11-23): MDP + heuristic joint relay (R-VE)
%   selection and time allocation in OFDM symbols per slot; R-VE dynamically chosen each slot among the
%   source and members that decoded; success per slot is an independent outage event (Eq. 33-34, no
%   combining); error-free HARQ feedback.
With Option~2, the source retransmits on its selected resources when feedback triggers it, and the HARQ
process context remains between the source and the receivers. Kuo's blind helper
scheme~\cite{kuo2021reliable} lets the source reserve up to two retransmission resources with its initial
transmission and assign them to a group member selected from the members' positions, which retransmits
the TB blindly if it decoded it, without group feedback. We assume
that the receivers combine the copies (feedback-less HARQ), but the source does not recover the
retransmission resources if the helper misses the packet. Leibel \emph{et al.}~\cite{leibel2020harq}
let the source use the HARQ feedback to select a member that decoded the packet and forward it to
out-of-range members, as a separate multi-hop forwarding step that is not combined with the source's
transmissions. Takahashi \emph{et al.}~\cite{aoki2026enhancing} let a relay retransmit a broadcast packet
on the source's reserved retransmission resources once an explicit ACK confirms that the relay decoded it,
assuming error-free ACKs; a fallback after relay failure is mentioned but not evaluated. Kim
\emph{et al.}~\cite{kim2019efficient} established the benefits of dynamically selecting the retransmitting
vehicle among the source and the members that decoded the packet, but model each retransmission as an
independent decoding trial without combining, with error-free feedback and an abstract OFDM-symbol-level
allocation. In contrast, our approach operates at
the 3GPP slot and TB level and preserves the sidelink HARQ process across the relay, so that the
destination members can soft-combine at the PHY.

Unlike standard 3GPP Rel-17/18 sidelink relaying, which operates above RLC via the sidelink relay
adaptation protocol (SRAP), our scheme works at the PHY and MAC layers. This preserves HARQ process
continuity across the relay and enables helper-agnostic, source-associated soft combining, early
termination on receiver feedback, and source recovery through fallback or takeover.
% NOTE: add 1-2 sentences on system-level NR V2X evaluations (5G-LENA) if a reviewer would expect them.

% =====================================================================================================
\section{System Model}\label{sec:model}
% NOTE: about 0.75 column + Fig. 1 (scenario sketch).
\subsection{Scenarios}
% NOTE: checked against the code: 9 vehicles (numPlatoon = 9), 35 m between members (5 m car + 30 m gap),
% highway KPI window 3-52.7 s (about 50 s). \todo{verify the platoon speed (~110 km/h) from the SUMO
% trace, and the 1 km background area on the highway (the code shows 500 m: mcr = 500)}
We use ns-3 with the 5G-LENA sidelink module~\cite{ali20213gpp} and SUMO~\cite{sumo} mobility traces. A
platoon of nine vehicles, each 5\,m long with 30\,m gaps, groupcasts from its leader; apart from the
helper, when one is used, the leader is the only platoon member that transmits. We define three
scenarios. On the \emph{highway}, with three lanes per direction, the platoon drives in the middle lane at
about \todo{110}\,km/h. In the two other scenarios (Fig.~\ref{fig:scenario}) the platoon crosses an urban
roundabout surrounded by buildings: it enters from the north and leaves to the west after a quarter of
the ring (\emph{Roundabout~W}) or to the east after three quarters (\emph{Roundabout~E}), so on
Roundabout~E the platoon stays bent around the buildings longer. On the roundabouts we measure from when
the first platoon member enters the curve until the last member leaves it; on the highway over a 50\,s
window.

\begin{figure}[t]
  \centering
  \fbox{\parbox[c][3.5cm]{0.9\columnwidth}{\centering \todo{scenario figure: highway and roundabout W/E
  with the platoon, buildings and background vehicles}}}
  \caption{Scenarios. \todo{caption}}
  \label{fig:scenario}
\end{figure}

\subsection{Traffic}
% NOTE: checked against the code: PDB 100 ms periodic, 50 ms aperiodic; 60 % of the background vehicles
% periodic (bgPeriodicFrac = 0.6); loads 0/15/30/45/60 vehicles; pool occupancy up to 0.69 (cbr_occupied).
% \todo{verify the background placement radius per scenario (code default mcr = 500 m)}
The leader's traffic follows TR~37.885~\cite{3gpp_tr37885}:
\begin{enumerate}
  \item \emph{Periodic}: a packet every 100\,ms (a 300-byte packet, then four of 190 bytes), with a
        100\,ms packet delay budget (PDB);
  \item \emph{Aperiodic}: a packet every 50\,ms plus an exponential time of mean 50\,ms, of 200 to
        1200 bytes in 200-byte steps, with a 50\,ms PDB.
\end{enumerate}
Every packet fits in one TB (about \todo{2} of the 5 subchannels on average for periodic traffic, almost
all 5 for aperiodic traffic), and RLC runs in unacknowledged mode (no segmentation, no ARQ), so a packet
is received exactly when its TB is decoded and the per-packet PRR equals the TB delivery ratio. The leader
uses semi-persistent reservations for periodic traffic and dynamic grants for aperiodic traffic; when a
packet does not fit the current reservation, the leader selects new resources. The pool is shared with
0, 15, 30, 45 or 60 background vehicles placed at random for each seed within \todo{1\,km} of the platoon
on the highway and \todo{500}\,m of the roundabout centre. Of these vehicles, 60\,\% send periodic traffic
on semi-persistent reservations and 40\,\% aperiodic traffic on dynamic grants. The share of
slot--subchannel cells they occupy (pool CBR) grows with their number, up to 69\,\% at 60 vehicles.

\subsection{Channel and Radio Configuration}
% NOTE: checked against the code (nr_sl_stack.cc, run.csv): 5.89 GHz, 20 MHz dedicated SL carrier,
% mu = 1, 5 subchannels of 10 PRBs, MCS 14, 23 dBm, 1x2 isotropic antenna at 1.5 m, sensing window
% 100 ms, T2 = 33 slots, sl-MaxNumTx = 3, PSFCH period 1 slot. The 5-slot minimum spacing of feedback
% retransmissions measured in the traces is consistent with a 3-slot PSFCH gap + 2-slot processing.
% \todo{verify: RSRP threshold -128 dBm, PSSCH-to-PSFCH gap 3 slots, processing gap 2 slots}
Path loss, shadowing and fast fading follow TR~37.885, with vehicles and buildings as
blockers~\cite{3gpp_tr37885}. The carrier is 5.89\,GHz with 20\,MHz, a 30\,kHz subcarrier spacing
(0.5\,ms slots) and 23\,dBm transmit power. We use a fixed modulation and coding scheme (MCS)~14, five
subchannels of ten resource blocks, HARQ incremental redundancy, and a selection window of 2 to 33
slots. The PSFCH period is one slot, with a three-slot PSSCH-to-PSFCH gap and a two-slot processing gap.
Table~\ref{tab:params} lists the settings.

\subsection{Performance Metrics}
% NOTE: define precisely, the reviewers will check:
%   group PRR      = mean over members of the share of packets received within the PDB;
%   worst-member PRR = PRR of the worst member in the run;
%   transmissions per packet = leader + helper PSSCH transmissions / packets;
%   pool CBR = share of (slot, subchannel) cells of the pool used by any UE in the KPI window
%              (each cell counted once; NOT the offered load, which can exceed 1).
\todo{Metric definitions}
\begin{equation}
  \mathrm{PRR}_{\text{group}} = \frac{1}{N}\sum_{m=1}^{N}\frac{|\{k : d_{m,k}\le T_{\text{PDB}}\}|}{K}
  \label{eq:prr}
\end{equation}

% =====================================================================================================
\section{Schemes Compared}\label{sec:schemes}
% NOTE: one subsection per scheme, or one paragraph each + Table I and a timing figure (Fig. 2) that shows
% the slots of one TB: leader copy, handed slots, PSFCH occasions, takeover / fallback.
\subsection{Source Option 2 HARQ}\label{sec:r3}
\todo{Baseline: the leader retransmits on any NACK; stops on all-ACK.}

\subsection{Blind Helper Relaying}\label{sec:r4c}
\todo{Helper selected per TB (Kuo's rule), named in SCI-2A; every reserved slot after the leader's
first copy is handed to the helper, which relays a decoded TB blindly; receivers soft-combine.}

\subsection{Blind Helper Relaying with Fallback}\label{sec:r4cfb}
\todo{After its own copy the leader listens in the helper's next handed slot; no helper SCI there means the
helper missed the TB, and the leader takes back the following handed slot. Spaced variant: slots selected
with the PSFCH minimum time gap so the takeover is always feasible.}
% NOTE: decide whether this is a full scheme in the paper or an ablation in Sec. VI-D.

\subsection{Proposed Cooperative HARQ}\label{sec:r6}
\todo{Feedback-driven helper relaying: helper relays on a member NACK; a helper NACK makes the leader take
the next handed slot back; a silent helper triggers the fallback; transmission stops on all-ACK.}
% NOTE: add pseudo-code (Algorithm 1) for the leader and helper decisions if space allows.

\begin{table}[t]
\centering
\caption{Schemes compared}\label{tab:schemes}
\footnotesize
\setlength{\tabcolsep}{3pt}
\begin{tabular}{@{}>{\raggedright\arraybackslash}p{2.5cm}cccc@{}}
\toprule
Scheme & Helper relays & Feedback & Takeover & Fallback \\
\midrule
\Scheme{R3}      & --      & Opt.~2 & --          & --         \\
\Scheme{R4+C}    & always  & --     & --          & --         \\
\Scheme{R4+C+fb} & always  & --     & --          & \checkmark \\
\Scheme{R6}      & on NACK & Opt.~2 & helper NACK & \checkmark \\
\bottomrule
\end{tabular}
\end{table}

% =====================================================================================================
\section{Simulation Setup}\label{sec:setup}
% NOTE: ns-3 version, 5G-LENA NR module version + what you changed (list the model extensions:
% helper naming in SCI-2A, slot handover, soft combining of relayed copies, PSFCH error modes, forced
% helper miss); SUMO scenarios (highway, roundabout W, roundabout E); seeds (10), run length, KPI window,
% confidence intervals (95 % Student-t over seeds; paired by seed where the schemes share the trace).
\todo{Simulator, scenarios, seeds, statistics}

\begin{table}[t]
\centering
\caption{Simulation parameters}\label{tab:params}
\begin{tabular*}{\columnwidth}{@{\extracolsep{\fill}}ll@{}}
\toprule
Parameter & Value \\
\midrule
Carrier, bandwidth & 5.89\,GHz, 20\,MHz (dedicated) \\
Numerology & 30\,kHz subcarriers, 0.5\,ms slots \\
Subchannels & 5 (10 PRBs each) \\
Transmit power, antenna & 23\,dBm; $1\times2$ isotropic, 1.5\,m high \\
MCS & 14, fixed \\
Sensing window, RSRP threshold & 100\,ms, $-128$\,dBm \todo{verify} \\
Selection window $[T_1,T_2]$ & [2, 33] slots \\
Transmissions per TB & up to 3 (1 + 2 retransmissions) \\
PSFCH period, min.\ gap & 1 slot, 3 slots \todo{verify} \\
HARQ & incremental redundancy \\
Platoon & 9 vehicles, 5\,m long, 30\,m gaps \\
Background vehicles & 0, 15, 30, 45, 60 (60\,\% periodic) \\
Pool CBR (measured) & 0.01--0.69 \\
Traffic, PDB & periodic 100\,ms; aperiodic 50\,ms \\
Seeds & 10 per point \\
\bottomrule
\end{tabular*}
\end{table}

% =====================================================================================================
\section{Results}\label{sec:results}
% NOTE: about 2.5 columns. Lead every subsection with the claim, then the figure, then the number.

\subsection{Reliability Against the Channel Load}
% NOTE: Fig. prr_vs_cbr (group PRR vs pool CBR, 3 routes x 2 traffic). Claims to consider:
%  - R4+C is the most load-sensitive scheme: its blind copies add collisions [aperiodic, highway,
%    0 -> 60 vehicles: 96.7 -> 68.3 %; on the roundabouts it falls to R3's level at CBR 0.70].
%  - R6 stays on top at every CBR [10-13 pp above R3 and R4+C on the roundabouts at CBR 0.70, aperiodic].
%  - periodic traffic suffers far less (announced reservations).
\todo{Group PRR vs CBR}
\begin{figure*}[t]
  \centering
  \includegraphics[width=0.9\textwidth]{prr_vs_cbr/prr_vs_cbr.pdf}
  \caption{Group PRR against the measured pool CBR. \todo{caption: what to look at}}
  \label{fig:prr_cbr}
\end{figure*}

\subsection{The Worst Member}
% NOTE: Fig. worst_vs_cbr. Claims to consider:
%  - R6 has the highest worst member everywhere [e.g. 30.9 / 49.6 % at 60 vehicles, aperiodic, W / E,
%    vs 23.9 / 39.2 % for R4+C and 15.1 / 22.3 % for R3].
%  - R3's worst member is limited by blockage, not by load, on the roundabouts.
\todo{Worst-member PRR}
\begin{figure}[t]
  \centering
  \includegraphics[width=\columnwidth]{prr_vs_cbr/worst_vs_cbr.pdf}
  \caption{Worst-member PRR against pool CBR. \todo{caption}}
  \label{fig:worst_cbr}
\end{figure}

\subsection{Cost: Transmissions per Packet}
% NOTE: table or bar plot of tx/pkt per scheme at 0 and 60 vehicles; R6 cheaper than the blind arms
% [R4+C+fb 2.6-2.9 tx/pkt vs R6 1.7-2.7]. Reliability per transmission is a good single number.
\todo{Resource cost}

\subsection{Where the Gain Comes From: the Leader Fallback}
% NOTE: ablation with R4+C+fb (paired by seed). Claims to consider:
%  - the fallback alone gives 61-93 % of R6's gain on the highway, 14-76 % on the roundabouts;
%  - the gain grows with the load [highway aperiodic +2.3 pp at 0 veh -> +10.6 pp at CBR 0.70];
%  - what is left to R6 on the roundabouts is the feedback-driven reach to blocked members.
% Figure: r4c_fallback/fallback_gain_vs_cbr.pdf (or a table if space is short).
\todo{Fallback ablation}

\subsection{Robustness to Helper Failure}
% NOTE: Fig. helper_miss_fb/prr_vs_p.pdf. The key distinction: payload loss (the helper NACKs) vs SCI
% loss (the helper is silent). Claims to consider:
%  - without feedback the two failure types are identical; R6 reacts to a NACK at once and barely
%    loses PRR [highway 99.8 -> 99.6 % at p = 0.7];
%  - with a silent helper, R6 and the fallback arm recover the same way (timeout-based takeover);
%  - resource selection matters: blind selection can place the handed slots back to back, leaving no
%    time to take one back [~10 % of fallbacks]; selecting with the PSFCH minimum time gap removes it,
%    after which the spaced blind arm matches R6 with a silent helper.
\todo{Helper failure}
\begin{figure}[t]
  \centering
  \includegraphics[width=\columnwidth]{helper_miss_fb/prr_vs_p.pdf}
  \caption{Group PRR against the forced helper miss probability $p$; solid: the helper fails the
  payload and NACKs, dashed: the helper misses the SCI and stays silent. \todo{caption}}
  \label{fig:helper_miss}
\end{figure}

\subsection{Robustness to PSFCH Errors}
% NOTE: optional if space allows: R6 reliability robust to lost / flipped feedback; cost rises
% [helper relays after all-ACK: 0.29 / 0.54 per packet at 10 / 30 % loss].
\todo{PSFCH errors (or move to a sentence in VI-E)}

% =====================================================================================================
\section{Discussion}\label{sec:discussion}
% NOTE: 3-4 short paragraphs: (1) when a helper pays off (blockage, load); (2) when it does not
% (open highway at low load: leader-only Option 2 is enough); (3) standardisation impact (fields in
% SCI-2, slot handover, min-gap rule); (4) limitations: one platoon, fixed helper rule, 3GPP CBR
% definition differs from the occupancy measure, simulated rather than measured channel.
\todo{Discussion and limitations}

% =====================================================================================================
\section{Conclusion}\label{sec:conclusion}
% NOTE: 1 paragraph: restate the problem, the protocol, the three headline numbers, one line of
% future work (analytical model, multiple platoons, helper's own resources).
\todo{Conclusion}

% -----------------------------------------------------------------------------------------------------
\section*{Acknowledgment}
\todo{Funding / acknowledgment, or delete}

\bibliographystyle{IEEEtran}
\bibliography{refs}

\end{document}
